% ============================================================
% Chapter 1 — Introduction (target 4–5 pp)
% Status: FULL DRAFT 2026-07-17, written fresh against the final
% campaign results (the 2026-07-06 chat draft was never recovered;
% if it resurfaces, merge selectively).
% ============================================================
\chapter{Introduction}
\label{ch:intro}

\section{Motivation}
\label{sec:motivation}

Reinforcement learning trains control policies by trial and error, and
in any non-trivial domain the designer faces a scheduling question
before the first gradient step: when the tasks to master span a range of
difficulty, \emph{in what order should they be presented}? Two answers
are natural. \emph{Curriculum learning} presents tasks from easy to
hard, echoing human pedagogy: build competence where success is
reachable, then extend it. \emph{Mixed training} presents everything
from the start --- the direct analogue of training on a shuffled
dataset, which is the unquestioned default in supervised learning. Both
strategies consume the same budget; choosing between them is, in
principle, free. If the choice changes what the learned policy does, it
is among the cheapest levers a designer controls --- and one of the
least characterized, especially in multi-agent settings, where most of
the curriculum literature's assumptions (single learner, sparse rewards,
stationary environment) do not hold.

Urban autonomous driving is a demanding place to ask the question.
It offers a natural difficulty axis (route length), heterogeneous
agents with different embodiments and constraints (vehicles and
pedestrians), a built-in tension between performance and safety
(faster policies complete more routes and hit more things), and a
canonical out-of-distribution test (drive on a town never seen in
training). It is also intrinsically multi-agent: policies do not just
share a world, they shape each other's environment while learning ---
the non-stationarity at the heart of multi-agent reinforcement
learning. This thesis asks the ordering question there, and treats the
training regime itself as the experimental variable.

\section{Research Question and Hypotheses}
\label{sec:rq}

\begin{quote}
  \emph{Does curriculum learning from easy to medium to hard produce
  measurably different multi-agent driving behaviour than mixed-batch
  training, at equal training budget?}
\end{quote}

The question is deliberately about \emph{behaviour}, not about a single
success-rate leaderboard: the comparison is evaluated along four axes
--- sample efficiency, final performance with its failure composition,
within-run training dynamics, and map generalization --- always
separately for vehicles and pedestrians. Three falsifiable hypotheses
structure the analysis:

\begin{itemize}
  \item \textbf{H1 (curriculum).} Staged exposure yields better early
        sample efficiency --- easy tasks provide reachable success to
        bootstrap from --- with a known risk of plateauing on, or
        forgetting, harder tasks introduced late.
  \item \textbf{H2 (mixed).} Uniform exposure covers the full task
        distribution from step zero, at the cost of a noisier early
        signal dominated by failures on hard tasks; its expected payoff
        is robustness --- what is seen during training does not shift.
  \item \textbf{H3 (interaction, secondary).} The size --- possibly the
        direction --- of the regime effect may depend on how value is
        estimated. In centralized-training architectures the critic is
        the component that digests multi-agent structure; replacing its
        encoder (flat MLP vs.\ graph-based message passing over agents)
        may change what task ordering does.
\end{itemize}

The campaign's short answers, developed in \cref{ch:results}: H1's
efficiency clause fails (all cells learn equally fast early --- dense
reward shaping already provides the signal curricula usually supply)
while ordering instead selects \emph{which} policy emerges, cautious
versus fast-and-colliding; H2's robustness clause fails (broader
coverage bought no transfer advantage); and H3 holds in strong form ---
under a GNN critic the in-domain regime effect vanishes entirely, and
only the generalization advantage of the curriculum survives across both
architectures.

\section{Approach at a Glance}
\label{sec:approach}

All experiments run in CARLA 0.9.16 on the urban map Town03, with
Town05 reserved for the generalization test. Six agents learn
simultaneously in one shared world: three vehicles and three
pedestrians, each assigned its own route per episode --- a design we
call \emph{parallel task coverage}, which multiplies experience per
episode and makes agent interaction a first-class phenomenon rather
than a nuisance. It is the fixed frame of both experimental arms, not a
variable. Policies are trained with MAPPO under centralized training
with decentralized execution: one shared policy per class, and a
centralized critic that sees all agents' observations during training
only. Difficulty is the target route length (30/60/\SI{100}{m});
success is strict route completion.

The two regimes share every component except the level-selection rule.
The \emph{curriculum} arm starts easy and unlocks harder levels through
a windowed competence gate (with budget-based pressure caps that
guarantee progression); the \emph{batch} arm interleaves all levels
uniformly from the first episode. The final campaign is a $2 \times 2$
factorial --- \{MLP, GNN critic\} $\times$ \{curriculum, batch\} ---
of four 3-million-step runs sharing one seed, paired across cells, each
followed by a 400-episode evaluation including Town05. With one run per
cell, every contrast is a seed-paired case study: effects are read
against a same-seed variance band measured on identical replicates, and
the primary contrast is observed twice, once per critic architecture.
The campaign sits on top of a longer gate-driven development phase ---
documented in full, including its failures and the environment bugs it
uncovered --- which is part of the thesis's methodological content
(\cref{ch:methodology}).

\section{Contributions}
\label{sec:contributions}

\begin{enumerate}
  \item \textbf{A controlled equal-budget comparison framework} for
        training regimes in multi-agent driving: seed-paired runs with
        byte-identical configurations up to the experimental flags,
        strict per-agent measurement rules, and a three-layer
        evaluation (training stream, static recomputation, final
        checkpoint evaluation with map transfer).
  \item \textbf{A gate-driven development methodology with an honest
        registry}: single-knob candidates promoted or reverted against
        predeclared quantitative gates, mechanism and outcome recorded
        separately, and environment-validity defects --- including one
        that moved the headline metric by $\approx 50$ percentage
        points with no policy change --- documented with their
        detection signatures.
  \item \textbf{An empirical characterization of what task ordering
        does} in this setting: no early-efficiency effect, but
        selection between qualitatively different driving equilibria;
        a curriculum advantage on the strict success criterion under
        the MLP critic; regime-insensitive pedestrians; and a
        platform-wide late-training erosion measured at fixed task
        level.
  \item \textbf{Evidence that regime effects are conditional on value
        estimation} (H3): swapping the critic encoder from MLP to GNN
        erases the in-domain curriculum advantage, rescues the batch
        cell, and leaves map generalization as the only component that
        replicates across architectures.
\end{enumerate}

\section{Thesis Outline}
\label{sec:outline}

\Cref{ch:background} introduces the necessary background: reinforcement
learning and PPO, multi-agent learning with centralized critics,
curriculum learning and its mixed-training control, and the CARLA
simulator. \Cref{ch:system} describes the experimental platform ---
environment, observation and action spaces, reward design, route
generation, the two training regimes, and the critic variants.
\Cref{ch:methodology} defines the measurement rules and the gate-driven
protocol, documents the environment-validity episodes, and specifies
the campaign design. \Cref{ch:results} reports the results along the
four comparison axes, closing with the verdicts on H1--H3.
\Cref{ch:discussion} interprets the findings and states threats to
validity and scope limits; \cref{ch:conclusions} concludes and lists
future work. Appendices provide the condensed experiment registry, full
configurations, per-run result tables, and reproducibility details.
